\section{The complete one-variable comparison algorithm}
We now supply the order decision that the normal-form proof deliberately
did not assume. There are two sorts in the auxiliary comparison language.
The source sort $P$ is $\mathcal P_1$. A range expression $R$ is a
nonempty finite word of kernels applied to one open $P$ expression.
Every argument of a kernel is either $P$ or $R$; a sum of range expressions
is not introduced into a kernel argument. Closed words are evaluated and
folded into constants. Every open expression of either sort is at least
$x$ on a computable tail. All generated masses and indices are fixed
natural or ordinal data on each branch.

\begin{theorem}[Effective eventual signs]\label{u:thm:sign}
Given two expressions in these sorts, one computes a sign
$\sigma\in\{-1,0,1\}$ and $N\ge1$ such that their difference has sign
$\sigma$ for every $x\ge N$. In particular this holds for every pair in
$\mathcal P_1$.
\end{theorem}
The proof is a simultaneous comparison and separation-certificate algorithm.
We specify its rules before proving their common termination measure.
Every rule retains the thresholds of all its children and the finitely
many kernel domains it uses. Since every open argument is at least $x$
eventually, each fixed genuine-domain obligation is met by increasing
$N$ to a computable constant.

\subsection{The comparison rules}
For $P/P$, form the signed normal-form row. An affine row is decided
directly. If tags remain, replace each relevant tag by its exact peak
word \eqref{u:display:4-5}, compare those $R$ expressions recursively and choose the
greatest peak. Distinct tags cannot have equal peaks beyond the bound of
Theorem~\ref{u:thm:zeros}. Return its coefficient sign, using
Lemma~\ref{u:lem:cells}. Identical forms return exact equality.

For $R/P$, a constant or nonconstant affine $P$ is eventually smaller:
the open $R$ is at least $f(x)$, which dominates every fixed affine expression.
For a nonlinear $P$ compare $R$ with its displayed peaks \eqref{u:display:4-5}, and use
the exact range-cell test \eqref{u:display:3-2}. The $P/R$ rule is symmetric.

For $R/R$ with equal outer principal, compare the two literal arguments
recursively. Strictly ordered arguments give the same strict kernel order
by Theorem~\ref{u:thm:adjacent}, or by monotonicity of $f$ at $\omega$.
If the arguments agree eventually, use the following same-input key:
mass is irrelevant at $\omega$; compare $\lceil c/2\rceil$ in reverse at
$\omega^2$; compare masses in reverse above $\omega^2$ by
Theorem~\ref{u:thm:mass}. Equal keys give equal outer values.

For different outer principals $H>M$, if $H$ has a successor exponent,
use the exact cut \eqref{u:eq:cut} on the $H$ side until the new outer
principal is $M$ or has a limit exponent still above $M$.
This cannot jump below $M$: every principal below $\theta\omega$ is at
most $\theta$. It takes only the finite terminal successor part of the
exponent. Combine these rewrites immediately with the equal-principal
rule or the following mixed-limit rule; do not recursively call the
unmodified comparison at an intermediate, larger syntax state.

In the mixed-limit case write
\begin{equation}
                  X=\KK Lc(A),\qquad Y=\KK Me(B),\quad M<L,
                                                               \label{u:display:8-1}
\end{equation}
where $L$ has a nonzero limit exponent. Recursively compare $X$ with $B$.
If $X\le B$, then $Y>B\ge X$. If $X>B$, the separation certificate
below, with lower profile $h_M$, gives
$X>h_M(B)\ge Y$, by Lemma~\ref{u:lem:direct}. There is no equality branch
in this mixed-limit case.

\subsection{The simultaneous limit-cut certificate}
Whenever one original side of a comparison is a bare limit-exponent
kernel $X=\KK Lc(A)$ and its result is $X>S$, the algorithm also
constructs, for every requested fixed lower profile $P$ below $L$, a
computable modulus $N_{X,S}(P)$ such that
\begin{equation}
                   X(x)>P(S(x))\quad(x\ge N_{X,S}(P)). \label{u:display:8-2}
\end{equation}
This concerns generated $P/R$ expressions $S$, not arbitrary integers
below $X$: the assertion for $S=X-1$ would be false, and subtraction
does not occur in the generated language.

Here is the certificate's complete case analysis, following the same
strictly smaller children as comparison.
\begin{enumerate}
\item If $S=ax+b$, use direct growth \eqref{u:display:5-3} for the fixed profile
$z\mapsto P(az+b)$ at the actual argument $A\ge x$. Positivity of $P$
makes this imply \eqref{u:display:8-2}.
\item If $S$ is nonlinear positive, $S<X$ forces every child peak $Z_i<X$
by the range-cell rule. Apply each smaller child's certificate to
$P\circ h_\omega$ and use $S<h_\omega(\max_i Z_i)$. Take the finite
maximum of their moduli; no permanent choice of maximizing child is needed.
\item If $S$ has the same outer principal $L$, strict order means either
a greater input for $X$, or equal inputs and smaller mass. The
lexicographic conclusion of Theorems~\ref{u:thm:adjacent} and
\ref{u:thm:limitmass} applies directly to $P$, at the genuine integer inputs.
\item If $S=\KK Me(B)$ has smaller outer principal, its comparison used
the smaller child $X>B$. Apply that child's certificate to
$P\circ h_M$ and use $S\le h_M(B)$. This handles any finite nesting of
lower contexts by recursive composition, not by an unproved absorption step.
\item If $S$ has larger limit outer principal $H>L$, its comparison with
$X$ compares $S$ against $A$. The only branch giving $X>S$ has $S\le A$;
direct growth for $\KK Lc$ gives $X>P(A)\ge P(S)$. In the other branch
the larger kernel's child certificate would force
$S>h_L(A)\ge X$, so cannot yield $X>S$.
\item A larger successor outer principal on $S$ is unfolded exactly until
it reaches $L$ or a larger limit exponent; use the corresponding preceding
case. The rewrites are exact, including their changed masses and domains.
\end{enumerate}
Every requested profile remains a fixed positive expression with heads
below the relevant limit $L$: only affine maps, fixed lower heads and
$h_\omega$ are composed. Thresholds are profile dependent, and are computed
by following these finitely many smaller children.

\subsection{A common well-founded measure}
Expand sharing when counting occurrences. Let $h$ be the number of
original source $h_\lambda$ nodes across the two compared expressions,
including their argument nodes. Let $\mathcal Q$ be the finite collection
of principal indices in the Cantor decompositions of all original source
heads and initial auxiliary kernels. Close it under finite successor-exponent
chains, stopping at exponent one or a limit exponent. No rule introduces
a principal outside this finite set. List it increasingly as
$\pi_0<\cdots<\pi_{m-1}$ and assign a kernel occurrence of index $\pi_i$
weight $\omega^{m-1-i}$. Let $W$ be their natural sum across the pair.
Use the lexicographic measure
\begin{equation}
                             (h,W).                 \label{u:display:8-3}
\end{equation}

Replacing an original tag by its peak removes that displayed $h$ node,
retaining its argument only once in each child comparison. In $P/P$ and
$R/P$ children, other siblings are not duplicated. Thus $h$ decreases.
Equal-principal comparison removes both outer kernels, decreasing $W$.
Mixed-limit comparison removes the lower outer kernel, also decreasing $W$.
A combined mixed-successor step removes the opposite lowest-rank $M$
occurrence and adds only finitely many occurrences of principals strictly
greater than $M$, retaining the original higher occurrence internally.
Their total weight is strictly smaller than one $M$ weight. Hence $W$
decreases. Affine or closed leaves stop.

Since $h$ is natural and $W<\omega^m$, every branch terminates. Branching
is finite, so the comparison tree is finite. The certificate uses that
same tree: each invocation in its case analysis belongs to a strictly
smaller comparison child or is a direct supplier application. Thus no
unchanged comparison state calls its own certificate. Mutual induction
on \eqref{u:display:8-3} proves both comparison correctness and \eqref{u:display:8-2}, establishing
Theorem~\ref{u:thm:sign}.

\subsection{Closure and exact finite exceptions}
The two sorts are closed under every recursive child. Peak replacement
introduces a finite kernel word on a positive argument plus its fixed
norm constant. Successor expansion introduces only a longer kernel word.
No signed inner offset, exposed count, division, inverse iterate or
variable ordinal index occurs. Increasing a mass changes fixed domain
constants without changing the finite rank set or measure.

\begin{example}[A nested comparison with the counts retained]
Consider $E=h_{2\omega}(x)$ and $F=h_\omega(h_\omega(x))$.
Deleting the two $\omega$ blocks gives
\[
 E=f(f(x+4))+2,\qquad F=f(f(x+2)+3)+1.
\]
Thus their peaks are respectively $f(f(x+4))$ and
$f(f(x+2)+3)$. The equal-outer-$\omega$ rule compares the arguments
$f(x+4)$ and $f(x+2)+3$. The latter positive nonlinear expression lies
strictly between its peak $f(x+2)$ and the next range value on the
computed tail. Comparing $f(x+4)$ with that peak reduces to the affine
comparison $x+4>x+2$. Hence $f(x+4)>f(x+2)+3$, and the original signed-row
rule returns $E>F$ on the effective tail. The offset $3$ has been paid
by a range-cell test, rather than deleted as negligible without proof.
The final procedure evaluates the finitely many earlier inputs exactly.
\end{example}

\begin{example}[Why the mass key changes at degree three]
At the same genuine input, the masses $1$ and $2$ give equal
$\omega^2$ kernels because their ceilings after division by two agree.
For a parent $\omega^3$ comparison with mass difference one, the first
common $\omega^2$ block can likewise tie when its retained child masses
are an odd $a$ and $a+1$. The next block reduces both child masses by
three; its ceiling classes are then distinct. The guaranteed two common
blocks ensure that those reduced masses remain nonnegative.
This is the two-block mechanism in
Theorem~\ref{u:thm:mass}, and explains why the algorithm uses neither strict
mass order at degree two nor a ceiling equivalence at all higher degrees.
\end{example}

\begin{corollary}[One-input quantified comparison]\label{u:cor:boolean}
For any finite Boolean combination of equality, weak-order and strict-order
atoms between expressions of $\mathcal P_1$, each of its existential and
universal closures over $x\in\NN$ is decidable.
\end{corollary}
\begin{proof}
Apply Theorem~\ref{u:thm:sign} to every atom and take the maximum tail bound.
The Boolean truth value is constant beyond that bound. Evaluate the
original positive input expressions exactly at every earlier positive
integer. For an existential statement, a true prefix instance or a true
tail suffices; for a universal statement require all prefix instances and
the tail. Evaluate original expressions, rather than auxiliary kernels
before their raised genuine domains.
\end{proof}

